\documentclass{article}
\usepackage{graphicx} % Required for inserting images
\usepackage{soul} % Required for Equation Strikethrough
\usepackage{amssymb}
\usepackage{amsmath} % Required for Equation Alignment
\usepackage{graphicx} % Required for Image
\title{Learning LATEX}
\author{Makhzum - Bin - Harun}
\date{November 2026}
\begin{document}
%Hello World!
\maketitle

\section{Section}
\subsection{Subsection}

This is my first LaTeX document.
I am learning LaTeX.\\
This is paragraph one.\\
This is paragraph two.

\section{Different Types of Text}
\subsection{Text Style}\\
\textbf{This text is bold.}\\
\textit{This text is Italic.}\\
\underline{This text is Underlined.}\\
\st{This text is Strikethrough.}
\subsection{Font Style}
{\tiny Tiny}\\
{\small Small}\\
{\large Large}\\
{\Large Larger}\\
{\Huge Huge}

\section{Different Types of Lists}
\subsection{Bullet List}

\begin{itemize} \item Apple \item Mango \item Orange \end{itemize}

\subsubsection{Bullet List Types}
\begin{itemize} \item[$\blacksquare$]Apple \item[$\square$]Mango \item Orange \end{itemize}

\subsection{Numbered List}
 \begin{enumerate} \item First \item Second \item Third \end{enumerate}

\subsection{Description List}
\begin{description}
  \item[Item 1] This is Item 1.
  \item[Item 2] This is Item 2.
\end{description}

\subsection{Nested List}
\begin{enumerate}
  \item Top-level ordered item
  \begin{itemize}
    \item Indented sub-bullet item
    \item Another sub-bullet item
  \end{itemize}
  \item Second top-level ordered item
\end{enumerate}

\section{Mathematics}
\subsection{Inline math \$...\$ or $ \(...\backslash\) $}
The equation is $a+b=c$.
Displayed math with slash \[a+b=c\]
\subsubsection{Power}
\paragraph{Superscript}
\[x^2\]
\paragraph{Subscript}
$x_1$
\subsubsection{Fraction}
$\frac{1}{2}$
\subsubsection{Square Root}
$\sqrt{x}$
\subsubsection{Summation}
$\sum_{i=1}^{n}$
\subsubsection{Integral}
$\int_0^1 x dx$
\subsection{Greek Letters}
\subsubsection{Small Greek Letters}
$\alpha$
$\beta$
$\gamma$
$\lambda$
$\theta$
$\pi$
$\sigma$
$\omega$
\subsubsection{Capital Greek Letters}
$A$
$B$
$\Gamma$
$\Lambda$
$\Theta$
$\Pi$
$\Sigma$
$\Omega$

\subsection{Equation Environment\textbackslash{begin\{equation\}}}
\begin{equation}
E=mc^2 \end{equation}
Automatically numbers equations.

\subsection{Align Equations}
\begin{align}
a &= b+c\\
x &= y+z
\end{align}

\section{Images}
\includegraphics[width=0.5\textwidth]{logo (1).png}

\section{Figure}
\begin{figure}
\centering
\includegraphics[width=.5\textwidth]{image}
\caption{My image}
\label{fig:image}
\end{figure}

\section{Tables}
\begin{tabular}{|c|c|}
\hline
Name & Age\\
\hline
John & 25\\
\hline
Alice & 30\\
\hline
\end{tabular}


\section{References}
Label
\label{sec:intro}
Reference
See Section~\ref{sec:intro}.
\section{Citations}
Using BibTeX
According to \cite{smith2023}.
\subsection{Bibliography}
\bibliography{references}
\section{Packages}
Packages add features. Examples:\\
\textbackslash{usepackage\{amsmath\}}\\
\textbackslash{usepackage\{graphicx\}}\\
\textbackslash{usepackage\{booktabs\}}\\
\textbackslash{usepackage\{hyperref\}}
\subsection{Common packages}
\begin{tabular}{|c|c|}
\hline
Package & Purpose\\
\hline
amsmath & Better mathematics\\
\hline
graphicx & Images\\
\hline
booktabs & Better tables\\
\hline
hyperref & Clickable links\\
\hline
xcolor & Colors\\
\hline
geometry & Page margins\\
\hline
float & Better figure placement\\
\hline
caption & Caption customization\\
\hline
\end{tabular}
\section {Special Characters}
Some symbols are reserved and must be escaped.
Character Latex Write # \# $ \$ % \% & \& _ \_ { \{ } \} ^ \^{} ~ \~{} \ \textbackslash

\end{document}



% This is ignored



19. Hyperlinks \usepackage{hyperref}
Then \href{https://example.com}{Visit Website}

20. Page Break \newpage
21. Line Break
Hello
\
\
World
22. Center \begin{center}
Hello \end{center}
23. Left and Right Alignment
Left \begin{flushleft}
Text \end{flushleft}
Right \begin{flushright}
Text \end{flushright}
24. Common Document Structure \documentclass{article} \usepackage{amsmath} \usepackage{graphicx} \usepackage{hyperref} \title{My Paper} \author{Your Name}

Page 6 of 9

\date{\today}
\begin{document}
\maketitle
\tableofcontents
\section{Introduction}
Text...
\section{Methods}
Text...
\section{Results}
Text...
\section{Conclusion}
Text...
\bibliographystyle{plain}
\bibliography{references}
\end{document}

25. The 20 Most Important Commands to Memorize
Command Purpose
\documentclass{} Document type
\usepackage{} Load a package
\begin{document} Start document
\end{document} End document
\section{} Section heading
\subsection{} Subsection
\textbf{} Bold
\textit{} Italic
\underline{} Underline
\item List item
\includegraphics{} Insert image
\caption{} Figure/table caption
\label{} Create a reference label
\ref{} Reference a label
\cite{} Cite a reference
\frac{}{} Fraction
\sqrt{} Square root
\sum Summation
\int Integral
\alpha, \beta, etc. Greek letters
